% Vietnamese translation for OI Public Library
% Source-PDF: IOI中国国家候选队论文/2009/汤可因/浅析竞赛中一类数学期望问题的解决方法.pdf
\documentclass[11pt]{article}
\usepackage{oipl-vn}

\newcommand{\E}{\operatorname{E}}
\newcommand{\Prob}{\operatorname{P}}
\newcommand{\Oh}{\mathcal{O}}

\title{Bàn sơ về phương pháp giải một lớp bài toán kỳ vọng trong thi đấu}
\author{Tang Keyin\\Trường Trung học số 8 Phúc Châu, Phúc Kiến}
\date{2009}

\begin{document}
\maketitle

\origtitle{Bàn sơ về phương pháp giải một lớp bài toán kỳ vọng trong thi đấu}
\sourcepath{IOI China candidate team papers / 2009 / Tang Keyin}

\begin{abstract}
Kỳ vọng có rất nhiều ứng dụng trong đời sống, và trong Olympic Tin học cũng
không ít bài yêu cầu tính kỳ vọng. Bài viết quy nạp các phương pháp thường
dùng cho bài toán kỳ vọng của biến ngẫu nhiên rời rạc thành hai nhóm lớn:
dùng truy hồi hoặc quy hoạch động, và lập hệ phương trình tuyến tính. Sau đó
phân tích, tổng kết qua nhiều ví dụ.
\end{abstract}

\paragraph{Từ khóa.}
Kỳ vọng; truy hồi; quy hoạch động; hệ phương trình.

\section{Mở đầu}

Trong xác suất và thống kê, kỳ vọng của một biến ngẫu nhiên rời rạc là tổng
các kết quả có thể xảy ra nhân với xác suất tương ứng. Kỳ vọng thường không
bằng bất kỳ kết quả cụ thể nào, nhưng là một đại lượng trực quan để đánh giá
một quá trình.

Vì kỳ vọng xuất hiện rộng rãi trong thực tế, các bài toán tính kỳ vọng cũng
xuất hiện nhiều trong thi tin học. Phần lớn các bài này là bài về biến ngẫu
nhiên rời rạc. Bài viết này tổng kết các phương pháp thường gặp cho lớp bài
đó.

\section{Kiến thức chuẩn bị}

\subsection{Định nghĩa}

Nếu \(X\) là biến ngẫu nhiên rời rạc, các giá trị có thể là
\(x_1,x_2,\ldots\), và xác suất tương ứng là \(p_1,p_2,\ldots\), với tổng xác
suất bằng \(1\), thì
\[
  \E(X)=\sum_i p_i x_i.
\]
Ví dụ, gieo một con xúc xắc đều, \(X\) là số điểm gieo được. Khi đó
\[
  \E(X)=1\cdot\frac16+2\cdot\frac16+\cdots+6\cdot\frac16=3.5.
\]

\subsection{Tính tuyến tính}

Với mọi biến ngẫu nhiên \(X,Y\) và hằng số \(a,b\),
\[
  \E(aX+bY)=a\E(X)+b\E(Y).
\]
Nếu \(X,Y\) độc lập và đều có kỳ vọng xác định, thì
\[
  \E(XY)=\E(X)\E(Y).
\]

\subsection{Công thức xác suất toàn phần}

Giả sử \(\{B_n\mid n=1,2,3,\ldots\}\) là một phân hoạch hữu hạn hoặc đếm được
của không gian xác suất. Với mọi sự kiện \(A\),
\[
  \Prob(A)=\sum_n \Prob(A\mid B_n)\Prob(B_n).
\]

\subsection{Kỳ vọng có điều kiện và công thức kỳ vọng toàn phần}

Ký hiệu
\[
  p_{ij}=\Prob(X=x_i,Y=y_j).
\]
Kỳ vọng có điều kiện của \(Y\) khi \(X=x_i\) được viết là
\(\E(Y\mid X=x_i)\). Khi đó
\[
\begin{aligned}
  \E(\E(Y\mid X))
  &=\sum_i \Prob(X=x_i)\E(Y\mid X=x_i)\\
  &=\sum_i\sum_k y_k p_{ik}\\
  &=\E(Y).
\end{aligned}
\]
Vì vậy
\[
  \E(Y)=\sum_i \Prob(X=x_i)\E(Y\mid X=x_i).
\]

Ví dụ, một công việc nếu chỉ do A làm thì trung bình cần \(4\) giờ; B có xác
suất \(0.4\) đến giúp, và khi hai người cùng làm thì trung bình cần \(3\) giờ.
Nếu \(X\) là số người làm việc và \(Y\) là thời gian hoàn thành, thì
\[
  \E(Y)=\Prob(X=1)\E(Y\mid X=1)+\Prob(X=2)\E(Y\mid X=2)
  =0.6\cdot4+0.4\cdot3=3.6.
\]

\section{Dùng truy hồi hoặc quy hoạch động}

Truy hồi là một thuật toán cơ bản. Với một phần bài toán kỳ vọng, ta không cần
liệt kê mọi khả năng mà có thể dùng kỳ vọng của các trạng thái đã biết để suy
ra trạng thái khác, hoặc dựa trên các trường hợp có kết quả giống nhau để tính
xác suất.

\subsection{Ví dụ 1: Congcong và Coco}

Trong một khu rừng ma thuật có mèo Congcong và chuột Coco. Khu rừng là một đồ
thị vô hướng gồm \(N\) cảnh điểm. Coco ở điểm \(M\). Mỗi đơn vị thời gian, Coco
chọn đều một trong các đỉnh kề hoặc đứng yên. Congcong ở điểm \(C\), luôn chọn
một đỉnh kề gần Coco hơn; nếu có nhiều lựa chọn, chọn đỉnh có số nhỏ nhất. Nếu
sau bước đầu chưa bắt được Coco, trong cùng đơn vị thời gian Congcong đi thêm
một bước nữa theo quy tắc đó. Mỗi đơn vị thời gian Congcong đi trước, Coco đi
sau. Hỏi trung bình sau bao nhiêu bước Congcong bắt được Coco.

\paragraph{Phân tích.}
Vì Congcong luôn đi theo đường ngắn nhất tới Coco và các cạnh không trọng số,
ta chạy BFS từ mỗi đỉnh để tiền xử lý \(p[i,j]\): nếu Congcong ở \(i\), Coco ở
\(j\), thì bước đầu tiên của Congcong đi tới đỉnh kề nào của \(i\) trên một
đường ngắn nhất tới \(j\), ưu tiên chỉ số nhỏ.

Đặt \(f[i,j]\) là số bước kỳ vọng để bắt được Coco khi Congcong ở \(i\), Coco
ở \(j\). Gọi \(w[j,k]\) là đỉnh kề thứ \(k\) của \(j\), và \(t[j]\) là bậc của
\(j\). Vị trí của Congcong sau hai bước là \(p[p[i,j],j]\). Coco có xác suất
\(1/(t[j]+1)\) đi tới mỗi đỉnh kề hoặc đứng yên. Vì vậy:
\[
  f[i,j]=
  \frac{
    \sum_{k=1}^{t[j]} f[p[p[i,j],j],w[j,k]]
    +f[p[p[i,j],j],j]
  }{t[j]+1}
  +1.
\]
Biên là \(f[i,i]=0\). Nếu \(p[i,j]=j\) hoặc \(p[p[i,j],j]=j\), Congcong bắt
được Coco trong bước hiện tại, nên \(f[i,j]=1\). Có thể dùng nhớ hóa để tính.

\paragraph{Tiểu kết.}
Với nhiều bài, có thể thiết kế trạng thái kỳ vọng và viết trực tiếp quan hệ
truy hồi. Thực chất đây là ứng dụng của công thức xác suất toàn phần và các
tính chất của kỳ vọng. Phần lập hệ phương trình tuyến tính bên dưới cũng bắt
nguồn từ ý tưởng này, nhưng xử lý thêm tình huống có chu trình.

\subsection{Ví dụ 2: Highlander}

Chọn ngẫu nhiên đều một hoán vị sai vị trí của \(1,2,\ldots,n\), tức không có
phần tử nào đứng yên. Hoán vị \(i_1i_2\cdots i_n\) tương ứng với đồ thị có
đỉnh \(1,\ldots,n\) và cạnh có hướng \((1,i_1),(2,i_2),\ldots,(n,i_n)\). Hỏi
kỳ vọng số đỉnh trên chu trình dài nhất. Giới hạn \(2\le n\le100\).

\paragraph{Phân tích.}
Đồ thị thu được có mọi đỉnh bậc vào và bậc ra đều bằng \(1\), nên nó gồm nhiều
chu trình rời nhau. Vì là hoán vị sai vị trí, không có vòng tự khuyên.

Ta xác định dần các chu trình theo độ dài tăng dần để tránh lặp. Nếu còn
\(m\) số chưa xác định và muốn chọn \(k\) số thuộc cùng một chu trình độ dài
\(k\), có \(A_m^k/k\) phương án. Nếu có \(i\) chu trình cùng độ dài, cần chia
thêm cho \(i!\).

Đặt \(f[i,j,k]\) là số phương án khi đã xác định \(i\) số, độ dài chu trình
lớn nhất hiện tại là \(j\), và có \(k\) chu trình độ dài \(j\). Đặt
\[
  g[i,j]=\sum_{k=1}^{\lfloor i/j\rfloor} f[i,j,k].
\]
Chuyển trạng thái của bài gốc có dạng:
\[
f[i,j,k]=
\begin{cases}
\dfrac{A_n^i}{i}, & i=j,\ i>1,\ k=1,\\[0.8em]
\dfrac{f[i-j,j,k-1]\cdot A_{n-i+j}^{j}}{j\,k},
  & 2\le j\le i-2,\ 1<k\le\lfloor i/j\rfloor,\\[0.8em]
\displaystyle\sum_{\ell=2}^{\min(j,i-j+1)}
  g[i-j,\ell]\frac{A_{n-i+j}^{j}}{j},
  & 2\le j\le i-2,\ k=1,\\[0.8em]
0, & \text{trường hợp khác.}
\end{cases}
\]
Đáp án là
\[
  \frac{\sum_{j=2}^{n}\sum_{k=1}^{n} j\,k\,f[n,j,k]}
       {\sum_{j=2}^{n}\sum_{k=1}^{n} f[n,j,k]}.
\]
Các giá trị \(f\) có thể rất lớn, nhưng vì chỉ dùng để tính xác suất, kiểu số
thực vẫn đủ cho yêu cầu bài.

\paragraph{Ghi chú.}
Mẫu số trong công thức đáp án chính là tổng số hoán vị sai vị trí của \(n\)
phần tử. Nếu ký hiệu số đó là \(d[n]\), ta cũng có truy hồi quen thuộc
\[
  d[n]=(n-1)(d[n-1]+d[n-2]),\qquad d[1]=0,\ d[2]=1.
\]
Có thể nén trạng thái xuống hai chiều để đạt độ phức tạp tốt hơn; cách trình
bày ở trên giữ dạng trực quan nhất của phép đếm.

\paragraph{Tiểu kết.}
Khi khó tìm quan hệ truy hồi trực tiếp giữa các kỳ vọng, có thể quay về định
nghĩa
\[
  \E(X)=\sum_i p_i x_i.
\]
Với loại bài này, trạng thái thường lưu xác suất hoặc số phương án; chỉ số
trạng thái trực tiếp hoặc gián tiếp biểu thị giá trị của biến ngẫu nhiên.

\subsection{Ví dụ 3: RedIsGood}

Trên bàn có \(R\) lá đỏ và \(B\) lá đen, được xáo ngẫu nhiên. Mỗi lần lật một
lá: lá đỏ nhận \(1\) đô-la, lá đen mất \(1\) đô-la. Người chơi có thể dừng bất
cứ lúc nào. Hỏi theo chiến lược tối ưu, kỳ vọng nhận được bao nhiêu tiền.

\paragraph{Phân tích.}
Đặt \(f[i,j]\) là kỳ vọng tiền nhận được khi còn \(i\) lá đỏ và \(j\) lá đen.
Có hai quyết định: dừng hoặc tiếp tục lật. Dừng cho kỳ vọng \(0\). Nếu lật,
xác suất lá đỏ và lá đen lần lượt là \(i/(i+j)\), \(j/(i+j)\), nên:
\[
f[i,j]=
\begin{cases}
\max\left\{0,\,
(f[i-1,j]+1)\frac{i}{i+j}
+(f[i,j-1]-1)\frac{j}{i+j}\right\},
  & i>0,\ j>0,\\[0.6em]
f[i-1,0]+1, & i>0,\ j=0,\\
0, & i=0.
\end{cases}
\]
Đáp án là \(f[R,B]\).

\paragraph{Tiểu kết.}
Với bài kỳ vọng có quyết định tối ưu, thường dùng quy hoạch động. Trạng thái
vẫn biểu diễn kỳ vọng, còn kỳ vọng thể hiện trạng thái tốt hay xấu. Trong ví
dụ này, nếu kỳ vọng của việc lật tiếp nhỏ hơn \(0\), tốt hơn là dừng lại.

\section{Lập hệ phương trình tuyến tính}

\subsection{Mô hình cơ bản}

Cho đồ thị có hướng \(G=(V,E)\), mỗi đỉnh \(i\) có trọng số \(W_i\). Với mỗi
cạnh \((u,v)\) có xác suất \(P_{u,v}>0\). Nếu tổng xác suất các cạnh ra từ
đỉnh \(i\) là
\[
  S_i=\sum_{(i,j)\in E} P_{i,j},
\]
thì tại đỉnh \(i\) có xác suất \(1-S_i\) dừng lại. Trọng số một đường đi
\(\langle v_1,v_2,\ldots\rangle\) là \(\sum_i W_{v_i}\). Hỏi bắt đầu từ đỉnh
\(s\), kỳ vọng tổng trọng số của đường đi đến khi dừng là bao nhiêu.

Nếu đồ thị không có chu trình, có thể truy hồi theo thứ tự topo. Nhưng với đồ
thị có chu trình, số đường đi có thể là vô hạn.

Đặt \(F_{i,j}\) là kỳ vọng tổng trọng số khi xuất phát từ \(i\) và đi nhiều
nhất \(j\) bước. Với \(j=0\), \(F_{i,0}=W_i\). Khi \(j>0\),
\[
  F_{i,j}
  =\sum_{(i,k)\in E}P_{i,k}(F_{k,j-1}+W_i)
  +\left(1-\sum_{(i,k)\in E}P_{i,k}\right)W_i
  =\sum_{(i,k)\in E}P_{i,k}F_{k,j-1}+W_i.
\]
Nếu \(F_{i,j}\) hội tụ khi \(j\to\infty\), ký hiệu giới hạn là \(F_i\). Khi
đó:
\[
  F_i=\sum_{(i,j)\in E} P_{i,j}F_j+W_i.
\]
Tất cả các đẳng thức này tạo thành một hệ phương trình tuyến tính. Giải hệ
bằng khử Gauss cho độ phức tạp ổn định; hơn nữa, từ việc hệ có nghiệm duy
nhất hay không, ta có thể đánh giá kỳ vọng có tồn tại hay không.

Như mọi hệ tuyến tính, hệ này vẫn có thể vô nghiệm hoặc có vô số nghiệm. Tuy
nhiên, việc toàn hệ không có nghiệm duy nhất không nhất thiết nói rằng ẩn cần
tìm cũng không duy nhất; ví dụ \(a=1,\ b=b+1\) vẫn xác định được \(a\). Vì
mỗi đỉnh ứng với một ẩn, trong thực tế chỉ cần giữ những đỉnh mà đỉnh xuất
phát \(s\) có thể đi tới. Nếu ẩn cần tìm vẫn không có nghiệm duy nhất, thông
thường có thể xem kỳ vọng là không tồn tại, nhưng vẫn phải xét điều kiện riêng
của từng bài, như ví dụ Mario bên dưới.

Nếu trọng số nằm trên cạnh, với cạnh \((i,j)\) có trọng số \(W_{i,j}\), phương
trình đổi thành:
\[
  F_i=\sum_{(i,j)\in E} P_{i,j}(F_j+W_{i,j}).
\]

Một chi tiết cài đặt hữu ích là đưa ẩn cần tìm \(F_s\) xuống cuối thứ tự khử,
nhờ vậy có thể tránh bước thế ngược trong một số bài.

\subsection{Ví dụ 4: First Knight}

Một hình chữ nhật được chia thành \(m\times n\) ô. Cho xác suất
\(P_{i,j}^{(k)}\) với \(1\le k\le4\), biểu thị xác suất từ ô \((i,j)\) đi tới
\((i+1,j)\), \((i,j+1)\), \((i-1,j)\), \((i,j-1)\). Kỵ sĩ bắt đầu ở
\((1,1)\), mỗi bước đi theo xác suất đã cho, độc lập với trước đó. Hỏi kỳ vọng
số bước để tới \((m,n)\). Giới hạn \(1\le m,n\le40\).

Mô hình trực tiếp cho phương trình:
\[
  E_{i,j}=
  P_{i,j}^{(1)}E_{i+1,j}
  +P_{i,j}^{(2)}E_{i,j+1}
  +P_{i,j}^{(3)}E_{i-1,j}
  +P_{i,j}^{(4)}E_{i,j-1}
  +1.
\]
Số ẩn là \(mn\); khử Gauss thô có độ phức tạp \(\Oh(m^3n^3)\), không chấp
nhận được khi có nhiều bộ dữ liệu.

\paragraph{Tối ưu trực quan.}
Ma trận mở rộng rất thưa: mỗi phương trình chỉ liên quan nhiều nhất bốn ẩn.
Khi hệ số ở cột trụ của một hàng bằng \(0\), hai hàng không cần khử. Ngoài ra
có thể điều chỉnh thứ tự khử: đặt \(E_{1,1}\) cuối cùng để khỏi thế ngược, và
khử trước các ô có \((i+j)\bmod2=1\). Các ô cùng màu bàn cờ này không trực tiếp
phụ thuộc nhau, nên khử chúng tiết kiệm nhiều thời gian.

\paragraph{Tối ưu sâu hơn.}
Quan sát theo từng hàng. Với \(i=1\), do xác suất đi ra ngoài bằng \(0\), các
phương trình của hàng đầu có thể khử để biểu diễn \(E_{1,j}\) bằng các
\(E_{2,1},E_{2,2},\ldots,E_{2,n}\) và hằng số:
\[
  E_{1,j}=a_1E_{2,1}+a_2E_{2,2}+\cdots+a_nE_{2,n}+c_{1,j}.
\]
Thay biểu thức này vào các phương trình của hàng \(2\), ta tiếp tục biểu diễn
hàng \(2\) theo hàng \(3\). Lặp lại cho tới hàng \(m\), rồi giải và thế ngược.
Mỗi lần chỉ khử \(n\) phương trình với không quá \(2n+1\) hạng liên quan, nên
độ phức tạp đạt khoảng \(\Oh(n^3m)\); có thể cài trực tiếp trong quá trình khử
Gauss bằng cách dùng thứ tự ngược từ \(E_{m,n}\) về \(E_{1,1}\).

\subsection{Ví dụ 5: Mario}

Mario sống trong mê cung \(N\times M\). Ô có thể là tiền, quái vật, ống, tường
hoặc ô trống. Vào ô tiền thì nhận số tiền ghi trên đó, tiền không biến mất.
Vào ô quái vật thì mất một mạng. Vào cửa ống thì lập tức bị chuyển tới cửa ra
duy nhất của ống; có thể có nhiều cửa vào cùng một cửa ra. Tường không đi
được. Mario bắt đầu ở một ô trống, có \(3\) mạng; mỗi bước chọn đều một trong
bốn ô kề có thể đi. Trò chơi kết thúc khi mạng bằng \(0\) hoặc khi không thể
nhận thêm tiền. Hỏi kỳ vọng số tiền nhận được; nếu vô hạn thì in \(-1\).
Giới hạn \(1\le N,M\le15\).

\paragraph{Phân tích.}
Với bài có \(3\) mạng, cách thường dùng là phân lớp đồ thị theo số mạng còn
lại: các lớp tương ứng với \(3,2,1\) mạng. Trừ tường và cửa vào ống, mỗi ô
tạo một đỉnh ở từng lớp. Nếu vào quái vật ở lớp \(k>1\), cạnh đi tới lớp
\(k-1\); ở lớp \(1\), quái vật là trạng thái kết thúc. Các ô khác nối trong
cùng lớp. Ô tiền có trọng số bằng số tiền, các ô khác có trọng số \(0\).

Điều kiện kết thúc thứ hai là không thể nhận thêm tiền. Ta BFS từ mỗi ô hoặc
trên đồ thị phù hợp để xác định các ô không thể tới một ô tiền dương nào; các
ô đó cũng được xem là trạng thái kết thúc và không có cạnh ra. Khi lập phương
trình, có thể không đưa chúng vào đồ thị, chỉ xem đóng góp sau đó là \(0\).

Sau khi xử lý như vậy, lập hệ phương trình như mô hình cơ bản. Nếu kỳ vọng của
đỉnh xuất phát không xác định, do các trọng số không âm và mọi đỉnh còn lại
đều có thể tới tiền, có thể kết luận kỳ vọng là vô hạn và in \(-1\).

\section{Tổng kết}

Bài viết bắt đầu từ các bài có thể giải bằng truy hồi và quy hoạch động, qua
đó rút ra cách thiết kế trạng thái kỳ vọng thường gặp. Với những trường hợp có
chu trình khiến truy hồi trực tiếp không hiệu quả hoặc không xác định, ta lập
hệ phương trình tuyến tính và giải bằng khử Gauss. Khi áp dụng vào bài cụ thể,
cần chú ý khai thác cấu trúc của đề để tối ưu thứ tự khử, độ thưa của ma trận,
hoặc cách xây dựng trạng thái.

Không thể nói bài viết bao phủ mọi phương pháp giải bài kỳ vọng trong thi tin
học. Điều quan trọng hơn là nắm chắc khái niệm kỳ vọng, thiết kế phương pháp
linh hoạt, quan sát đặc điểm riêng của đề bài và tìm ra bản chất.

\section*{Tài liệu tham khảo}

\begin{enumerate}
  \item Wikipedia, \url{http://www.wikipedia.org/}.
  \item Ge Rong, \textit{Nghiên cứu về các bài toán xác suất và kỳ vọng}, báo
  cáo nghiên cứu đội tuyển quốc gia 2004.
\end{enumerate}

\end{document}
